\section{A monomial bound containing the prescribed powers}
\label{sec:monomial}

We use the commuting forms of Theorem~\ref{thm:forms} to construct a
monomial ideal containing the prescribed powers.  The construction
preserves the Hilbert function and increases every graded Betti number.
The first step explains why the original resolution remains minimal
over the polynomial ring acting through these forms.

\begin{lemma}[Transfer of a minimal resolution]
\label{mono:transfer}
Let $S=\C[x_1,\ldots,x_n]$, let $k$ be a field containing $\C$, and
let $\Delta$ be a finite-dimensional central division $k$-algebra.
Put $m=\dim_k\Delta$ and
$N=\Delta[x_1,\ldots,x_n]$, where the $x_i$ are central.
Suppose that $t_1,\ldots,t_n\in N_1$ commute and that
\begin{equation}
          N=\bigoplus_{\alpha\in\Z_{\geq0}^n}\Delta t^\alpha
 \label{mono:basis}
\end{equation}
is a graded direct-sum decomposition.  Give
$T=k[y_1,\ldots,y_n]$ a right action on $N$ by multiplication,
$u\cdot y_i=ut_i$.
For every homogeneous ideal $I\subseteq S$, the right $T$-module
$Q=N/IN$ satisfies
\begin{align}
 \dim_k Q_d&=m\dim_\C(S/I)_d
          &&(d\in\Z),\label{mono:dimension-transfer}\\
 \beta^T_{p,j}(Q)&=m\beta^S_{p,j}(S/I)
          &&(p\geq0,\ j\in\Z).
          \label{mono:betti-transfer}
\end{align}
Here right and left modules over the commutative ring $T$ are
identified when computing $\Tor$.
\end{lemma}

\begin{proof}
Commutativity of the $t_i$ and centrality of $k$ make the stated right
$T$-action well defined.  The map
\begin{equation}
 \Delta\otimes_kT\longrightarrow N,
 \qquad b\otimes y^\alpha\longmapsto bt^\alpha
 \label{mono:free-map}
\end{equation}
is a graded vector-space isomorphism by \eqref{mono:basis}.  It is
right $T$-linear because $(bt^\alpha)t_i=bt^{\alpha+e_i}$.
Thus $N$ is right $T$-free of rank $m$, with a basis in degree zero,
and
\begin{equation}
                         N_{>0}=NT_{>0}.
 \label{mono:positive-degree}
\end{equation}

The left $S$-action commutes with the right $T$-action by associativity
of multiplication in $N$.  Moreover,
$N=\Delta\otimes_\C S$ as a left $S$-module, so $N$ is left $S$-free
and therefore flat.  Each element of $I$ is central in $N$, whence
$IN$ is a two-sided ideal.  In particular it is a right $T$-submodule,
and
\[
               Q\simeq\Delta\otimes_\C(S/I).
\]
Taking a complex basis of $(S/I)_d$ proves
\eqref{mono:dimension-transfer}.

Let $F_\bullet\to S/I$ be a minimal graded $S$-free resolution.
Since $S$ is commutative, regard its terms as right $S$-modules and
tensor with the $(S,T)$-bimodule $N$.  Flatness gives a resolution
\[
                    F_\bullet\otimes_SN\longrightarrow Q
\]
by graded free right $T$-modules.  An entry $f$ of an original
differential acts on $N$ by left multiplication.  This map is right
$T$-linear, again by associativity.

Choose a $k$-basis $b_1,\ldots,b_m$ of $\Delta$, which is a right
$T$-basis of $N$ under \eqref{mono:free-map}.  If $f$ is nonzero and
has degree $e>0$, then each $fb_u$ belongs to $N_e$.  By the graded
decomposition \eqref{mono:basis}, its coefficients in the basis
$b_1,\ldots,b_m$ are homogeneous polynomials in $T$ of degree $e$.
They therefore lie in $T_{>0}$.  Every nonzero original differential
entry has positive degree, so the transferred resolution is minimal.
Each summand $S(-s)$ of $F_p$ becomes $m$ summands $T(-s)$.
Lemma~\ref{filt:minimal} proves \eqref{mono:betti-transfer} for every
$p,j$.  If $I=S$, both quotients and their resolutions may be taken
to be zero, and the same formulas hold.
\end{proof}

The next proposition supplies the monomial ideal used by the rest of
the proof.  Its ideal is allowed to depend on $I$, but it is fixed
simultaneously for all homological and internal degrees.

\begin{proposition}[Monomial reduction]
\label{mono:bound}
Let $S=\C[x_1,\ldots,x_n]$, let
$2\leq a_1\leq\cdots\leq a_n$, and let $I\subseteq S$ be a
homogeneous ideal containing a homogeneous regular sequence
$f_1,\ldots,f_n$ with $\deg f_h=a_h$.  Put
$P=(x_1^{a_1},\ldots,x_n^{a_n})$.
There is a monomial ideal $M\subseteq S$ containing $P$ such that
\begin{align}
 \dim_\C(S/M)_d&=\dim_\C(S/I)_d
                   &&(d\in\Z),\label{mono:hilbert-bound}\\
 \beta^S_{p,j}(S/I)&\leq\beta^S_{p,j}(S/M)
                   &&(p\geq0,\ j\in\Z).
                   \label{mono:betti-bound}
\end{align}
\end{proposition}

\begin{proof}
If $I=S$, take $M=S$.  Assume henceforth that $I$ is proper.
Apply Theorem~\ref{thm:forms} to the given regular sequence.  It
supplies $k$, $\Delta$, and commuting degree-one forms $t_h$ satisfying
the basis property \eqref{mono:basis} and homogeneous identities
\begin{equation}
 t_h^{a_h}=c_hf_h+\sum_{\ell<h}B_{h\ell}f_\ell
                       +\sum_{\ell>h}C_{h\ell}t_\ell,
 \qquad c_h\in k^\times,
 \label{mono:triangular}
\end{equation}
where the multipliers belong to $N=\Delta[x_1,\ldots,x_n]$.
In particular $C_{h\ell}$ is homogeneous of degree $a_h-1$ when
nonzero.  Define $T=k[y_1,\ldots,y_n]$, $m=\dim_k\Delta$,
$W=IN$, and $Q=N/W$ as in Lemma~\ref{mono:transfer}.

Under \eqref{mono:free-map}, view $N$ as the free $T$-module with
coefficient space $V=\Delta$.  In each total degree, order exponents
by placing the lexicographically smallest tuple
\[
                         (\alpha_n,\ldots,\alpha_1)
\]
first.  This order is compatible with addition.  Apply
Proposition~\ref{filt:comparison} to the homogeneous $T$-submodule
$W$ with this order.

For an exponent $\alpha$, let $V_\alpha\subseteq\Delta$ be its
leading coefficient space.  It is closed under left multiplication
by $\Delta$: if $w\in W$ witnesses a coefficient $b$ at $\alpha$
with no earlier coefficients, then $\lambda w\in W$ witnesses
$\lambda b$.  Here $W$ is left $\Delta$-linear because it is an ideal
of $N$.  Hence $V_\alpha$ is a left ideal of $\Delta$.  Since
$\Delta$ is a division algebra,
\[
                         V_\alpha=0\quad\hbox{or}\quad\Delta.
\]
Explicitly, if $0\ne b\in V_\alpha$, every $c\in\Delta$ is
$(cb^{-1})b$ and therefore also belongs to $V_\alpha$.

The set
\[
          \mathcal U=\{\alpha\in\Z_{\geq0}^n:V_\alpha=\Delta\}
\]
is upward closed under addition of exponent vectors, since the leading
object is a $T$-submodule.  Define $M\subseteq S$ to be the monomial
ideal with exponent set $\mathcal U$, and define
$M_k\subseteq T$ by the same exponent set.  Then
\begin{equation}
 N/\operatorname{in}_{\prec}W
       \simeq\Delta\otimes_k(T/M_k)
       \simeq(T/M_k)^{\oplus m}
 \label{mono:initial-quotient}
\end{equation}
as graded right $T$-modules.

We next verify the power containment.  Since every $f_\ell$ lies
in $I$, Equation~\eqref{mono:triangular} gives
\[
             t_h^{a_h}-\sum_{\ell>h}C_{h\ell}t_\ell\in W.
\]
Expand each $C_{h\ell}$ in the left $\Delta$-basis $t^\alpha$.
Right multiplication by $t_\ell$ turns a term $bt^\alpha$ into
$bt^{\alpha+e_\ell}$, using only commutativity of the forms.
Every resulting exponent has a positive coordinate of index greater
than $h$.  It consequently lies strictly after $a_he_h$ in the
chosen order, because $a_he_h$ has zero coordinates at all indices
greater than $h$.  The displayed element therefore has leading
position $a_he_h$ with coefficient $1$.  Thus $a_he_h\in\mathcal U$
for every $h$, proving $P\subseteq M$.

It remains to compare the numerical invariants.  The scalar extension
of $S/M$ from $\C$ to $k$, with $x_i$ renamed $y_i$, is $T/M_k$.
Its Hilbert function is unchanged.  Tensoring the Koszul complex over
$\C$ with the field $k$ also preserves exactness and commutes with
homology, so
\begin{align*}
 \dim_k(T/M_k)_d&=\dim_\C(S/M)_d,\\
 \beta^T_{p,j}(T/M_k)&=\beta^S_{p,j}(S/M).
\end{align*}
Equations \eqref{mono:dimension-transfer} and
\eqref{mono:initial-quotient}, together with
\eqref{filt:hilbert}, now give
\[
 m\dim_\C(S/I)_d
   =\dim_k Q_d
   =\dim_k(N/\operatorname{in}_{\prec}W)_d
   =m\dim_\C(S/M)_d.
\]
Likewise, \eqref{mono:betti-transfer},
\eqref{filt:betti}, and \eqref{mono:initial-quotient} yield, for
every $p\geq0$ and every $j\in\Z$,
\[
 m\beta^S_{p,j}(S/I)
  =\beta^T_{p,j}(Q)
  \leq\beta^T_{p,j}(N/\operatorname{in}_{\prec}W)
  =m\beta^S_{p,j}(S/M).
\]
Cancelling the positive integer $m$ proves both assertions.
\end{proof}

The right action in Lemma~\ref{mono:transfer} and the placement of the
$\Delta$-coefficients on the left are deliberate.  The argument uses
commutativity among the forms $t_i$, but never requires an individual
coefficient in $\Delta$ to commute with a form.  Proposition~\ref{mono:bound}
completes the reduction to an ideal containing $P$. The remaining
comparison is the prescribed-powers theorem of Mermin and Murai
\cite[Theorem~3.1]{MM}. Sections~\ref{sec:stabilization}--\ref{sec:tor}
give a direct proof by stabilization, box compression, and a Koszul-rank
induction.
